\subsection{The inflation is an information bound}\label{sec:information}

Proposition~\ref{prop:coupling} is a statement about one diagnostic, built on a
moment estimator. Its content is more general, and the generality matters for
whether the boundary is a property of an implementation or of the estimation
problem.

Subtracting sampling variance from between-population dispersion is exactly
what estimating the model variance of a Fay--Herriot small-area model does
\citep{fay1979estimates,rao2015small}. For area $c$, a direct estimate
$y_c=\theta_c+e_c$ with known sampling variance $D_c$, and
$\theta_c=x_c'\beta+u_c$ with $u_c\sim N(0,A)$. The Fisher information for $A$
with the $D_c$ known is $I(A)=\tfrac12\sum_c(A+D_c)^{-2}$, so any unbiased
estimator obeys
\begin{equation}\label{eq:crlb}
\mathrm{RSE}(\widehat A)\ \ge\ \frac{1}{A}\sqrt{\frac{2}{\sum_c(A+D_c)^{-2}}}
\ \ \overset{D_c\equiv D}{=}\ \ \sqrt{\frac{2}{K}}\cdot\frac{1}{1-\rho^2},
\qquad \rho^2=\frac{D}{A+D},
\end{equation}
the same $(1-\rho^2)^{-1}$ inflation, now from the information rather than from
a particular estimator. Requiring relative precision $\tau$ gives
$K\ge2/[\tau^2(1-\rho^2)^2]$, which is~\eqref{eq:boundary} up to the unit
arising from the moment estimator's $K-1$.

\paragraph{Attribution.} None of this is new, and it is important to say so
precisely. Equation~\eqref{eq:crlb} is the standard asymptotic variance of the
model-variance estimator in this model \citep{rao2015small}. The underlying
phenomenon---subtracting a component leaves the standard error unchanged while
shrinking what it is divided by, so relative precision degrades by the shrinkage
factor---is the familiar difficulty of variance-component estimation with a
small between-group component, and it is the same difficulty that makes the
between-study variance imprecise in random-effects meta-analysis
\citep{dersimonian1986meta}. \citet{partlett2017random} document the resulting
coupled condition empirically in that setting: intervals behave acceptably only
when the number of studies is large \emph{and} the heterogeneity is not
small relative to within-study error, the two requirements trading off exactly
as~\eqref{eq:boundary} describes.

What we add is not the algebra. It is three things: the transfer to a
\emph{calibration} setting, where the coupling governs whether a procedure may
be used at all rather than how wide an interval is; the closed form
\eqref{eq:boundary}, which turns the trade-off into a planning quantity
evaluated at the design share actually present; and the observation, developed
in Section~\ref{sec:applications}, that the feasibility criterion accompanying
the implementation under study is a fixed population count that
is~\eqref{eq:boundary} evaluated where the correction is unnecessary. The
implementation is stated in full in Section~\ref{sec:coupling}, and the
authors' relationship to it is described on the title page.

Two properties of~\eqref{eq:crlb} matter for practice, and we verify both on
\InfCells{} simulated conditions at \InfReps{} replicates
(Table~\ref{tab:crlb}).

\emph{The bound is attained by the estimator practitioners use.} Restricted
maximum likelihood tracks it closely: over conditions with $K\ge100$ the ratio
of realised to bound relative standard error has median \InfRatioMed{} and
range \InfRatioLo--\InfRatioHi. The inflation therefore cannot be evaded by a
better estimator, and the earlier restriction to unbiased estimation, which
excluded exactly the shrinkage estimators small-area practice is built on, is
not what carries the result.

\emph{The closed form is the equal-variance case, and it is conservative.}
With $D_c\equiv D$ it reproduces the exact bound to \InfClosedErr{} relative
error. With sampling variances dispersed across areas at the same mean, the
exact bound falls below it, by Jensen's inequality on $\sum_c(A+D_c)^{-2}$; the
closed form then overstates the requirement, so using it for planning is
conservative rather than optimistic.

\paragraph{A familiar symptom.} If the boundary is real, the regime it calls
infeasible should be recognisable to practitioners, and it is. The share of
replicates on which REML returns the boundary estimate $\widehat A=0$ is
\ZeroLoRho\% or less at design shares up to $0.5$ at every $K$ examined, and at
$\rho=0.9$ it is \ZeroHiRhoLoK\% at $K=30$, \ZeroHiRhoMidK\% at $K=100$ and
\ZeroHiRhoHiK\% at $K=500$. Zero and negative estimates of a between-area
variance are a long-standing nuisance in small-area work. They arise, on this
reading, in precisely the regime where the component is too small relative to
sampling error for the available number of areas to determine it, which is what
the boundary states.

\begin{table}[htbp]\centering
\caption{Relative standard error of the Fay--Herriot model variance:
information bound, REML, and the share of replicates returning a zero estimate.
Equal sampling variances, \InfReps{} replicates per condition.}\label{tab:crlb}
\begin{tabular}{rrrrrr}\toprule
& \multicolumn{2}{c}{$\rho=0.5$} & \multicolumn{3}{c}{$\rho=0.9$}\\
\cmidrule(lr){2-3}\cmidrule(lr){4-6}
$K$ & bound & REML & bound & REML & at zero\\\midrule
30 & 0.344 & 0.346 & 1.359 & 1.191 & \ZeroHiRhoLoK\%\\
100 & 0.189 & 0.183 & 0.744 & 0.708 & \ZeroHiRhoMidK\%\\
250 & 0.119 & 0.121 & 0.471 & 0.467 & 1.3\%\\
500 & 0.084 & 0.085 & 0.333 & 0.334 & \ZeroHiRhoHiK\%\\
\bottomrule\end{tabular}
\end{table}
